\chapter{萬章章句下}
\kaishi*

\section{論伯夷、伊尹、柳下惠、孔子}

孟子曰：「伯夷，目不視惡色，耳不聽惡聲。非其君不事，非其民不使。治則進，亂則退。橫政之所出，橫民之所止，不忍居也。思與鄉人處，如以朝衣朝冠坐於塗炭也。當紂之時，居北海之濱，以待天下之清也。故聞伯夷之風者，頑夫廉，懦夫有立志。」\jiazhu[1]{孟子反覆嗟伯夷、伊尹、柳下惠之德，以爲足以配於聖人，故數章陳之，猶詩人有所誦述至於數四，蓋其留意者也。義見上篇矣。此復言不視惡色，謂行不正而有美色者，若夏姬之比也。耳不聽惡聲，謂鄭聲也。後世聞其風者，頑貪之夫更思廉潔，懦弱之人更思有立義之志也。}

「伊尹曰：『何事非君？何使非民？』治亦進，亂亦進。曰：『天之生斯民也，使先知覺後知，使先覺覺後覺。予，天民之先覺者也。予將以此道覺此民也。』思天下之民，匹夫匹婦有不與被堯舜之澤者，若己推而內之溝中。其自任以天下之重也。」\jiazhu[1]{說與上同。}

「柳下惠不羞汙君，不辭小官。進不隱賢，必以其道。遺佚而不怨，阨窮而不閔。與鄉人處，由由然不忍去也。『爾爲爾，我爲我，雖袒裼裸裎於我側，爾焉能浼我哉？』故聞柳下惠之風者，鄙夫寬，薄夫敦。」\jiazhu[1]{鄙狹者更寬優，薄淺者更深厚。}

「孔子之去齊，接淅而行；去魯，曰：『遲遲吾行也，去父母國之道也。』可以速而速，可以久而久，可以處而處，可以仕而仕，孔子也。」\jiazhu[1]{淅，漬米也。不及炊，避惡亟也。魯，父母之國，遲遲不忍去也，是其道也。孔子聖人，故能量時宜，動中權也。}

孟子曰：「伯夷，聖之清者也；伊尹，聖之任者也；柳下惠，聖之和者也；孔子，聖之時者也。孔子之謂集大成。集大成也者，金聲而玉振之也。金聲也者，始條理也；玉振之也者，終條理也。」\jiazhu[1]{伯夷清，伊尹任，柳下惠和，皆得聖人之道也。孔子時行則行，時止則止。孔子集先聖之大道，以成己之聖德者也，故能金聲而玉振之。振，揚也。故如金聲之有殺，振揚玉音，始終如一也。始條理者，金從革，可治之使條理；終條理者，玉終其聲而不細也，含五德而不撓也。}

「始條理者，智之事也；終條理者，聖之事也。」\jiazhu[1]{智者知理物，聖人終始同。}

「智，譬則巧也；聖，譬則力也。由射於百步之外也，其至，爾力也；其中，非爾力也。」\jiazhu[1]{以智譬由人之有技巧也，可學而益之；以聖譬由力之有多少，自有極限，不可強增。聖人受天性，可庶幾而不可及也。夫射遠而至，爾努力也；其中的者，爾之巧也。思改其手，用巧意乃能中也。}\par \jiazhu[1]{章指言：聖人由力，力有常也；賢者由巧，巧可增也。仲尼天高，故不可階；他人丘陵，丘陵由可踰。所謂小同而大異者也。}

\section{北宫錡問周室班爵禄}

北宮錡問曰：「周室班爵禄也，如之何？」\jiazhu[1]{北宮錡，衛人。班，列也。問周家班列爵禄等差謂何？}

孟子曰：「其詳不可得聞也。諸侯惡其害己也，而皆去其籍。然而軻也嘗聞其略也。」\jiazhu[1]{詳，悉也，不可得備知也。諸侯欲恣行，憎惡其法度妨害己之所爲，故滅去典籍。今周禮司禄之官無其職，是則諸侯皆去之，故使不復存也。軻，孟子名。略，粗也。言嘗聞其大綱如此矣。今考之禮記王制，則合也。}

「天子一位，公一位，侯一位，伯一位，子、男同一位，凡五等也。」\jiazhu[1]{公謂上公九命及二王後也。自天子以下列尊卑之位，凡五等。}

「君一位，卿一位，大夫一位，上士一位，中士一位，下士一位，凡六等。」\jiazhu[1]{諸侯法天子，臣名亦有此六等，從君下至於士也。}

「天子之制，地方千里。公、侯皆方百里，伯七十里，子、男五十里，凡四等。不能五十里，不達於天子，附於諸侯，曰附庸。」\jiazhu[1]{凡此四等，土地之等差也。天子封畿千里，諸侯方百里，象雷震也。小者不能特達於天子，因大國以名通，曰附庸也。}

「天子之卿受地視侯，大夫受地視伯，元士受地視子、男。」\jiazhu[1]{視，比也。天子之卿、大夫、士所受采地之制也。}

「大國地方百里，君十卿禄，卿禄四大夫，大夫倍上士，上士倍中士，中士倍下士，下士與庶人在官者同禄，禄足以代其耕也。」\jiazhu[1]{公、侯之國爲大國。卿禄居於君禄十分之一也。上士之禄居大夫禄二分之一也。中士、下士轉相倍。庶人在官者，未命爲士者也，其禄比上農夫。士不得耕，以禄代耕也。}

「次國地方七十里，君十卿禄，卿禄三大夫，大夫倍上士，上士倍中士，中士倍下士，下士與庶人在官者同禄，禄足以代其耕也。」\jiazhu[1]{伯爲次國。大夫禄居卿禄三分之一也。}

「小國地方五十里，君十卿禄，卿禄二大夫，大夫倍上士，上士倍中士，中士倍下士，下士與庶人在官者同禄，禄足以代其耕也。」\jiazhu[1]{子、男爲小國。大夫禄居卿禄二分之一也。}

「耕者之所獲：一夫百畝。百畝之糞，上農夫食九人，上次食八人，中食七人，中次食六人，下食五人。庶人在官者，其禄以是爲差。」\jiazhu[1]{獲，得也。一夫一婦佃田百畝。百畝之田，加之以糞，是爲上農夫，其所得穀足以食九口。庶人在官者，食禄之等差由農夫有上、中、下之次，亦有此五等，若今之斗食、佐史、除吏也。}\par \jiazhu[1]{章指言：聖人制禄，上下差敘，貴有常尊，賤有等威。諸侯僭越，滅籍從私。孟子略記，言其大綱，以答北宮子之問。}

\section{萬章問友}

萬章問曰：「敢問友。」\jiazhu[1]{問朋友之道也。}

孟子曰：「不挾長，不挾貴，不挾兄弟而友。友也者，友其德也，不可以有挾也。」\jiazhu[1]{長，年長；貴，貴勢；兄弟，兄弟有富貴者。不挾是，乃爲友，謂相友以德也。}

「孟獻子，百乘之家也，有友五人焉：樂正裘、牧仲，其三人則予忘之矣。獻子之與此五人者友也，無獻子之家者也。此五人者，亦有獻子之家，則不與之友矣。」\jiazhu[1]{獻子，魯卿孟氏也，有百乘之賦。樂正裘、牧仲，其五人者，皆賢人無位者也。此五人者自有獻子之家富貴，而復有德，不肯與獻子友也。獻子以其富貴下此五人，五人屈禮而就也。}

「非惟百乘之家爲然也，雖小國之君亦有之。費惠公曰：『吾於子思，則師之矣；吾於顔般，則友之矣；王順、長息，則事我者也。』」\jiazhu[1]{小國之君，若費惠公者也。王順、長息，德不能見師友，故曰事我者也。}

「非惟小國之君爲然也，雖大國之君亦有之。晉平公於亥唐也，入云則入，坐云則坐，食云則食。雖疏食菜羹，未嘗不飽，蓋不敢不飽也。然終於此而已矣。」\jiazhu[1]{大國之君，如晉平公者也。亥唐，晉賢人也，隱居陋巷者。平公常往造之，亥唐言入，平公乃入；言坐，乃坐；言食，乃食也。疏食，糲食也。不敢不飽，敬賢也。終於此，平公但以此禮下之而已。}

「弗與共天位也，弗與治天職也，弗與食天禄也。士之尊賢者也，非王公尊賢也。」\jiazhu[1]{位、職、禄皆天之所以授賢者，而平公不與亥唐共之，而但卑身下之，是乃匹夫尊賢者之禮耳。王公尊賢，當與共天職矣。}

「舜尚見帝，帝館甥於貳室，亦饗舜，迭爲賓主，是天子而友匹夫也。」\jiazhu[1]{尚，上也。舜在畎畝之時，堯友禮之。舜上見堯，堯舍之於貳室。貳室，副宫也。堯亦就饗舜之所設，更迭爲賓主。禮謂妻父曰外舅，謂我舅者，吾謂之甥。堯以女妻舜，故謂舜甥。卒與之天位，是天子之友匹夫也。}

「用下敬上，謂之貴貴；用上敬下，謂之尊賢。貴貴、尊賢，其義一也。」\jiazhu[1]{下敬上，臣恭於君也；上敬下，君禮於臣也。皆禮所尚，故云其義一也。}\par \jiazhu[1]{章指言：匹夫友賢，下之以德；王公友賢，授之以爵。大聖之行，千載爲法者也。}

\section{萬章問交際}

萬章曰：「敢問交際何心也？」\jiazhu[1]{際，接也。問交接道當執何心爲可者。}

孟子曰：「恭也。」\jiazhu[1]{當執恭敬爲心。}

曰：「『卻之卻之爲不恭』，何哉？」\jiazhu[1]{萬章問：卻不受尊者禮，謂之不恭，何然也？}

曰：「尊者賜之，曰：『其所取之者，義乎？不義乎？』而後受之，以是爲不恭，故弗卻也。」\jiazhu[1]{孟子曰：今尊者賜己，己問其所取此物寧以義乎？得無不義，乃後受之，以是爲不恭，故不當問尊者不義而卻之也。}

曰：「請無以辭卻之，以心卻之，曰：『其取諸民之不義也。』而以他辭無受，不可乎？」\jiazhu[1]{萬章曰：請無正以不義之辭卻也，心知其不義，以他辭讓無受之，不可邪？}

曰：「其交也以道，其接也以禮，斯孔子受之矣。」\jiazhu[1]{孟子言其來求交己以道理，其接待己有禮者，若斯孔子受之矣。言可受也。}

萬章曰：「今有禦人於國門之外者，其交也以道，其餽也以禮，斯可受禦與？」\jiazhu[1]{禦人以兵禦人而奪之貨。如是而以禮道來接己，斯可受乎？}

曰：「不可。《康誥》曰：『殺越人于貨，閔不畏死，凡民罔不譈。』是不待教而誅者也。殷受夏，周受殷，所不辭也，於今爲烈。如之何其受之？」\jiazhu[1]{孟子曰：不可受也。《康誥》，尚書篇名，周公戒成王、康叔。封越、于皆於也。殺於人，取於貨，閔然不知畏死者。譈，殺也。凡民無不得殺之者也。若此之惡，不待君之教命，遭人得討之。三代相傳以此法，不須辭問也。於今爲烈，烈，明法。如之何受其餽也？}

曰：「今之諸侯取之於民也，猶禦也。苟善其禮際矣，斯君子受之，敢問何說也？」\jiazhu[1]{萬章曰：今諸侯賦稅於民，不由其道，履畝強求，猶禦人也。欲善其禮以接君子，君子欲受之，何說也？君子謂孟子。}

曰：「子以爲有王者作，將比今之諸侯而誅之乎？其教之不改而後誅之乎？夫謂非其有而取之者盜也，充類至義之盡也。孔子之仕於魯也，魯人獵較，孔子亦獵較。獵較猶可，而況受其賜乎？」\jiazhu[1]{孟子謂萬章曰：子以爲後如有聖人興作，將比地盡誅今之諸侯乎？將教之，其不改者乃誅之乎？言必教之，誅其不改者也。殷之衰亦猶周之末，武王不盡誅殷之諸侯，滅國五十而已。知後王者亦不盡誅也。謂非其有而竊取之者爲盜。充滿至甚也，滿其類，大過至者，但義盡耳，未爲盜也。諸侯本當稅民之類者，今大盡耳，亦不可比於禦。孔子隨魯人之獵較。獵較者，田獵相較奪禽獸，得之以祭，時俗所尚，以爲吉祥。孔子不違而從之，所以小同於世也。獵較尚猶可爲，況受其賜而不可也？}

曰：「然則孔子之仕也，非事道與？」\jiazhu[1]{萬章問：孔子之仕，非欲事行其道與？}

曰：「事道也。」\jiazhu[1]{孟子曰：孔子所仕者，欲事行其道。}

「事道奚獵較也？」\jiazhu[1]{萬章曰：孔子欲仕道，如何可獵較也？}

曰：「孔子先簿正祭器，不以四方之食供簿正。」\jiazhu[1]{孟子曰：孔子仕於衰世，不可卒暴改戾，故以漸正之。先爲簿書以正其宗廟祭祀之器，即其舊禮，取備於國中，不以四方珍食供其所簿正之器度。珍食難常有，乏絕則爲不敬，故獵較以祭也。}

曰：「奚不去也？」\jiazhu[1]{萬章曰：孔子不得行道，何爲不去？}

曰：「爲之兆也。兆足以行矣而不行，而後去。是以未嘗有所終三年淹也。」\jiazhu[1]{兆，始也。孔子每仕，常爲之正本造始，欲以次治之，而不見用，占其事始而退。足以行之矣而君不行也，然後則孔子去矣。終者，竟也。孔子未嘗得竟事一國也。三年淹留而不去者也。}

「孔子有見行可之仕，有際可之仕，有公養之仕。於季桓子，見行可之仕也；於衛靈公，際可之仕也；於衛孝公，公養之仕也。」\jiazhu[1]{行可，冀可行道也。魯卿季桓子秉國之政，孔子仕之，冀可得因之行道也。際，接也。衛靈公接遇孔子以禮，故見之也。衛孝公以國君養賢者之禮養孔子，故宿留以答之矣。}\par \jiazhu[1]{章指言：聖人憂民，樂行其道。苟善辭命，不忍逆距；不合則去，亦不淹久。蓋仲尼行止之節也。}

\section{孟子論仕}

孟子曰：「仕非爲貧也，而有時乎爲貧；娶妻非爲養也，而有時乎爲養。」\jiazhu[1]{仕本爲行道濟民也，而有以居貧親老而仕者。娶妻本爲繼嗣也，而有以親執釜竈，不擇妻而娶者。}

「爲貧者，辭尊居卑，辭富居貧。」\jiazhu[1]{爲貧之仕，當讓高顯之位，無求重禄。}

「辭尊居卑，辭富居貧，惡乎宜乎？抱關擊柝。」\jiazhu[1]{辭尊貧者，安所宜乎？宜居抱關擊柝監門之職也。柝，門關之木也，擊，椎之也。或曰：柝，行夜所擊木也。《傳》曰：「魯擊柝聞於邾。」}

「孔子嘗爲委吏矣，曰：『會計當而已矣。』嘗爲乘田矣，曰：『牛羊茁壯長而已矣。』位卑而言高，罪也；立乎人之本朝而道不行，恥也。」\jiazhu[1]{孔子嘗以貧而禄仕。委吏，主委積倉庾之吏也，不失會計，當直其多少而已。乘田，苑囿之吏也，主六畜之芻牧者也，牛羊茁壯肥好長大而已。茁，茁生長貌也。《詩》云：「彼茁者葭。」位卑不得高言豫朝事，故但稱職而已。立本朝，大道當行，不行爲己之恥。是以君子禄仕者，不處大位。}\par \jiazhu[1]{章指言：國有道，則能者處卿相；國無道，則聖人居乘田。量時安卑，不受言責，獨善其身之道也。}

\section{萬章問士不託諸侯}

萬章曰：「士之不託諸侯，何也？」\jiazhu[1]{託，寄也。謂若寄公，食禄於所託之國也。}

孟子曰：「不敢也。諸侯失國而後託於諸侯，禮也；士之託於諸侯，非禮也。」\jiazhu[1]{謂士位輕，本非諸侯敵體，故不敢比失國諸侯得爲寄公也。}

萬章曰：「君餽之粟，則受之乎？」\jiazhu[1]{士窮而無禄，君餽之粟，則可受之乎？}

曰：「受之。」\jiazhu[1]{孟子曰：受之也。}

「受之何義也？」\jiazhu[1]{萬章曰：受粟何義也？}

曰：「君之於氓也，固周之。」\jiazhu[1]{氓，民也。孟子曰：君之於民，固當周其窮乏，況於士乎？}

曰：「周之則受，賜之則不受，何也？」\jiazhu[1]{萬章言：士窮，君周之則受，賜之則不受，何也？周者，謂周急稟貧民之常科也；賜者，謂禮賜橫加也。}

曰：「不敢也。」\jiazhu[1]{孟子曰：士不敢受賜。}

曰：「敢問其不敢何也？」\jiazhu[1]{萬章問：何爲不敢？}

曰：「抱關擊柝者，皆有常職以食於上。無常職而賜於上者，以爲不恭也。」\jiazhu[1]{孟子曰：有職事者可食於上。禄士不仕，自以不任職事而空受賜，爲不恭，故不受也。}

曰：「君餽之則受之，不識可常繼乎？」\jiazhu[1]{萬章曰：君禮餽賢臣，賢臣受之，不知可繼續而常來致之乎？將當輒更以君命將之也。}

曰：「繆公之於子思也，亟問，亟餽鼎肉。子思不悅。於卒也，摽使者出諸大門之外，北面稽首再拜而不受，曰：『今而後知君之犬馬畜伋。』蓋自是臺無餽也。」\jiazhu[1]{孟子曰：魯繆公時尊禮子思，數問，數餽鼎肉。子思以君命煩，故不悅也。於卒者，末後復來時也。摽，麾也。麾使者出大門之外，再拜叩頭不受，曰：「今而後知君犬馬畜伋。」伋，子思名也。責君之不優以不煩，而但數與之食物若養犬馬。臺，賤官，主使令者。《傳》曰：「僕臣臺。」從是之後，臺不持餽來。繆公愠也。愠，恨也。}

「悅賢不能舉，又不能養也，可謂悅賢乎？」\jiazhu[1]{孟子譏繆公之雖欲有悅賢之意，而不能舉用使行其道，又不能優養終竟之，豈可謂能悅賢也？}

曰：「敢問國君欲養君子，如何斯可謂養矣？」\jiazhu[1]{萬章問國君養賢之法也。}

曰：「以君命將之，再拜稽首而受。其後廩人繼粟，庖人繼肉，不以君命將之。子思以爲鼎肉使己僕僕爾亟拜也，非養君子之道也。」\jiazhu[1]{將者，行也。孟子曰：始以君命行禮，拜受之。其後倉廩之吏繼其粟，將盡復送；厨宰之人日送其肉，不復以君命者，欲使賢者不答以敬，所以優之也。子思所以非繆公者，以爲鼎肉使己數拜故也。僕僕，煩猥貌，謂其不得養君子之道也。}

「堯之於舜也，使其子九男事之，二女女焉，百官、牛羊、倉廩備，以養舜於畎畝之中，後舉而加諸上位。故曰，王公之尊賢者也。」\jiazhu[1]{堯之於舜如是，是王公尊賢之道也。九男以下，已說於上篇。上位，尊帝位也。}\par \jiazhu[1]{章指言：知賢之道，舉之爲上，養之爲次。不舉不養，賢惡肯歸？是以孟子上陳堯舜之大法，下刺繆公之不弘也。}

\section{萬章問不見諸侯}

萬章曰：「敢問不見諸侯，何義也？」\jiazhu[1]{問諸侯聘請而夫子不見之，於義何取也？}

孟子曰：「在國曰市井之臣，在野曰草莽之臣，皆謂庶人。庶人不傳質爲臣，不敢見於諸侯，禮也。」\jiazhu[1]{在國，謂都邑也，民會於市，故曰市井之人；在野，野居之人，莽亦草也。庶，衆也。衆庶之人，未得爲臣。傳，執也。見君之質，執雉之屬也。未爲臣，則不敢見之禮也。}

萬章曰：「庶人，召之役，則往役；君欲見之，召之，則不往見之，何也？」\jiazhu[1]{庶人召使給役事，則往供事；君召之見，不肯往見，何也？}

曰：「往役，義也；往見，不義也。且君之欲見之也，何爲也哉？」\jiazhu[1]{孟子曰：庶人法當給役，故往役，義也。庶人非臣也，不當見君，故往見，不義也。且君何爲欲見之而召之也？}

曰：「爲其多聞也，爲其賢也。」\jiazhu[1]{萬章曰：君以是欲見之也。}

曰：「爲其多聞也，則天子不召師，而況諸侯乎？爲其賢也，則吾未聞欲見賢而召之也。」\jiazhu[1]{孟子曰：安有召師、召賢之禮而可往見也？}

「繆公亟見於子思，曰：『古千乘之國以友士，何如？』子思不悅，曰：『古之人有言曰：事之云乎，豈曰友之云乎？』子思之不悅也，豈不曰：『以位，則子君也，我臣也，何敢與君友也？以德，則子事我者也，奚可以與我友？』千乘之君求與之友而不可得也，而況可召與？」\jiazhu[1]{魯繆公欲友子思，子思不悅而稱曰：古人曰見賢人當事之，豈云友之邪？孟子云：子思所以不悅者，豈不謂臣不可友君，弟子不可友師也？若子思之意，亦不可友，況乎可召之？}

「齊景公田，招虞人以旌，不至，將殺之。志士不忘在溝壑，勇士不忘喪其元。孔子奚取焉？取非其招不往也。」\jiazhu[1]{已說於上篇。}

曰：「敢問招虞人何以？」\jiazhu[1]{萬章問招虞人當何用也。}

曰：「以皮冠。庶人以旃，士以旂，大夫以旌。」\jiazhu[1]{孟子曰：招禮若是。皮冠，弁也。旃，通帛也，因章曰旃。旂，旌有鈴者。旌，注旄首者。}

「以大夫之招招虞人，虞人死不敢往；以士之招招庶人，庶人豈敢往哉？況乎以不賢人之招招賢人乎？」\jiazhu[1]{以貴者之招招賤人，賤人尚不敢往，況以不賢人之招招賢人乎？不賢之招，不以禮也。}

「欲見賢人而不以其道，猶欲其入而閉之門也。夫義，路也；禮，門也。惟君子能由是路，出入是門也。」\jiazhu[1]{欲人之入而閉其門，何得而入乎？閉門，由閉禮也。}

「《詩》云：『周道如厎，其直如矢。君子所履，小人所視。』」\jiazhu[1]{《詩·小雅·大東》之篇。厎，平；矢，直；視，比也。周道平直，君子履直道，小人比而則之，以喻虞人能效君子守死善道也。}

萬章曰：「孔子，君命召，不俟駕而行。然則孔子非與？」\jiazhu[1]{俟，待也。孔子不待駕而應君命也。孔子爲之非與？}

曰：「孔子當仕有官職，而以其官召之也。」\jiazhu[1]{孟子言孔子所以不待駕者，孔子當仕位，有當職之事，君以其官名召之，豈得不顛倒？《詩》云：「顛之倒之，自公召之。」不謂賢者無位而君欲召見也。}\par \jiazhu[1]{章指言：君子之志，志於行道。不得其禮，亦不苟往。於禮之可，伊尹三聘而後就湯；道之未洽，沮、溺耦耕，接輿佯狂，豈可見也？}

\section{孟子論尚友}

孟子謂萬章曰：「一鄉之善士斯友一鄉之善士，一國之善士斯友一國之善士，天下之善士斯友天下之善士。」\jiazhu[1]{鄉，一鄉之善者；國，國中之善者；天下，四海之內也。各以大小來相友，自爲疇匹也。}

「以友天下之善士爲未足，又尚論古之人。頌其詩，讀其書，不知其人，可乎？是以論其世也。是尚友也。」\jiazhu[1]{好善者以天下之善士爲未足，極其善道也。尚，上也。乃復上論古之人。頌其詩，詩歌頌之，故曰頌。讀其書，猶恐未知古人高下，故論其世以別之也。在三皇之世爲上，在五帝之世爲次，在三王之世爲下。是爲好上友之人也。}\par \jiazhu[1]{章指言：好高慕遠，君子之道。雖各有倫，樂其崇茂。是以仲尼曰：「無友不如己者。」高山仰止，景行行止。}

\section{齊宣王問卿}

齊宣王問卿。孟子曰：「王何卿之問也？」\jiazhu[1]{王問何卿也。}

王曰：「卿不同乎？」曰：「不同。有貴戚之卿，有異姓之卿。」\jiazhu[1]{孟子曰：卿不同。貴戚之卿，謂內外親族也；異姓之卿，謂有德命爲三卿也。}

王曰：「請問貴戚之卿。」\jiazhu[1]{問貴戚之卿如何。}

曰：「君有大過則諫，反覆之而不聽，則易位。」\jiazhu[1]{孟子曰：貴戚之卿反覆諫君，君不聽，則欲易君之位，更立親戚之貴者。}

王勃然變乎色。\jiazhu[1]{王聞此言，愠怒而驚懼，故勃然變色。}

曰：「王勿異也。王問臣，臣不敢不以正對。」\jiazhu[1]{孟子曰：王勿怪也。王問臣，臣不敢不以其正義對。}

王色定，然後請問異姓之卿。\jiazhu[1]{王意解，顏色定，復問異姓之卿如之何。}

曰：「君有過則諫，反覆之而不聽，則去。」\jiazhu[1]{孟子言異姓之卿諫君不從，三而待放，遂不聽之，則去而之他國也。}\par \jiazhu[1]{章指言：國須賢臣，必擇忠良，親近貴戚，或遭殃禍。伊發有莘，爲殷興道。故云成湯立賢無方也。}
